% =============================================================================
% MATHEMATICAL ANALYSIS 06: Validity of the Dipole Approximation
% Stand-alone derivation for inspection
% =============================================================================

\documentclass[12pt]{article}
\usepackage{amsmath,amssymb,amsthm}
\usepackage[margin=1in]{geometry}
\usepackage{hyperref}

\newcommand{\Scfac}{a}
\newcommand{\Wight}{W}
\newcommand{\Nnum}{\mathcal{N}}
\newcommand{\pd}{\partial}
\newcommand{\be}{\begin{equation}}
\newcommand{\ee}{\end{equation}}
\newcommand{\lb}{\left}
\newcommand{\rb}{\right}

\title{Analysis 06: Validity of the Dipole Approximation}
\author{Calculation Notes --- Satishchandran \& Rao}
\date{}

\begin{document}
\maketitle

\section*{The dipole approximation}

The decoherence calculation uses the dipole approximation, replacing the current difference:
%
\begin{equation}
  j_{1}(x) - j_{2}(x)
  \approx
  q\,d(\tau)\,s^{a}\,\nabla_{a}\delta^{(3)}(\mathbf{x}-\mathbf{X}(\tau)) \,.
  \label{eq:dipole}
\end{equation}
%
This is the leading term in a multipole expansion in powers of the separation $d/\lambda$,
where $\lambda$ is the characteristic wavelength of field modes contributing to $\Gamma$.

\section*{Structure of the multipole expansion}

The exact current difference, for two point charges separated by $d\,\mathbf{s}$
along a unit vector $\mathbf{s}$, is:
%
\begin{align}
  j_{1}(\mathbf{x},\tau) - j_{2}(\mathbf{x},\tau)
  &=
  q\,\delta^{(3)}\!\lb(\mathbf{x}-\mathbf{X}-\tfrac{d}{2}\mathbf{s}\rb)
  - q\,\delta^{(3)}\!\lb(\mathbf{x}-\mathbf{X}+\tfrac{d}{2}\mathbf{s}\rb) \nonumber\\
  &=
  q\,d\,s^{i}\,\pd_{i}\delta^{(3)}(\mathbf{x}-\mathbf{X})
  + \frac{q}{24}d^{3}(s^{i}\pd_{i})^{3}\delta^{(3)}(\mathbf{x}-\mathbf{X})
  + \cdots
\end{align}
%
The expansion involves only odd powers of $d$ (by symmetry under $\mathbf{s}\to-\mathbf{s}$).
The first correction (octupole) is suppressed by $(d\,k)^{2}$ relative to the dipole,
where $k$ is the wavenumber.

\section*{Relevant field mode wavelengths}

The decoherence integral is dominated by field modes with conformal wavenumber
$k \sim 1/\Delta t$ (characteristic scale of the conformal-time integration domain).
In physical (proper) units, the corresponding wavelength is
%
\begin{equation}
  \lambda_{\rm phys} \sim \Scfac(\tau_{\rm on})\cdot\Delta t \,,
\end{equation}
%
where $\Delta t$ is the conformal time duration of the superposition.

\subsection*{Power-law case ($p>1$)}

The conformal time duration is bounded above by the horizon conformal time:
%
\begin{equation}
  \Delta t \leq t_{\infty} = \frac{u_{\rm on}^{1-p}}{C_{p}(p-1)} \,.
\end{equation}
%
The physical horizon distance at $\tau_{\rm on}$ is
$d_{H} = \Scfac(\tau_{\rm on})\cdot t_{\infty} = C_{p}u_{\rm on}^{p}\cdot t_{\infty}
= u_{\rm on}/(p-1)$ (proper distance to horizon).
The Hubble radius is $R_{H} = \Scfac/\dot{\Scfac} = (\tau_{\rm on}-\tau_{*})/p = u_{\rm on}/p$.

So the effective minimum wavelength (IR cutoff from the horizon) is:
%
\begin{equation}
  \lambda_{\rm eff} \sim d_{H} \sim \frac{u_{\rm on}}{p-1} \sim R_{H} \,.
\end{equation}
%
The dipole approximation requires:
%
\begin{equation}
  d_{0} \ll \lambda_{\rm eff} \sim R_{H} = \frac{u_{\rm on}}{p} \,.
  \label{eq:dipole-cond-PL}
\end{equation}

\subsection*{De Sitter case}

The Hubble radius is $R_{H} = 1/H$ (constant), and the effective mode wavelength is
also $\sim 1/H$.  The dipole condition is:
%
\begin{equation}
  d_{0} \ll \frac{1}{H} \,.
  \label{eq:dipole-cond-dS}
\end{equation}

\section*{Error estimate}

The next correction to the decoherence exponent from the octupole term scales as:
%
\begin{equation}
  \frac{\delta\Gamma_{\rm octupole}}{\Gamma_{\rm dipole}}
  \sim
  (d_{0}\,k_{\rm eff})^{2}
  \sim
  \lb(\frac{d_{0}}{R_{H}}\rb)^{2} \,.
\end{equation}
%
For a laboratory superposition of a mesoscopic object:
\begin{itemize}
\item $d_{0} \sim 10^{-9}$~m (nanometer separation)
\item $R_{H} \sim c/H_{0} \sim 10^{26}$~m (Hubble radius today)
\end{itemize}
The relative correction is $(d_{0}/R_{H})^{2} \sim 10^{-70}$: completely negligible.

\section*{The 1/3 factor from angular averaging}

In an isotropic FLRW spacetime, the spatial slices are homogeneous and isotropic
(maximally symmetric).  The decoherence exponent involves
$\hat{s}^{i}\hat{s}^{j}\langle\pd_{i}\pd_{j}\Wight\rangle$, where the angular average
is over the possible orientations of the dipole $\hat\mathbf{s}$.

For an isotropic kernel $K_{ij} = f(\tau_{1},\tau_{2})\delta_{ij}$:
%
\begin{equation}
  \hat{s}^{i}\hat{s}^{j}K_{ij}
  = f(\tau_{1},\tau_{2})\,|\hat\mathbf{s}|^{2} = f(\tau_{1},\tau_{2}) \,.
\end{equation}

However, the relevant tensor at coincident spatial points is not $\delta_{ij}$ but rather
comes from $\pd_{i}\pd_{j'}(1/\tilde\sigma)|_{x=x'} = 2\delta_{ij}/(\eta_1-\eta_2)^4$.
This is already proportional to $\delta_{ij}$, so
$\hat{s}^{i}\hat{s}^{j}\cdot(2\delta_{ij}/(\eta_1-\eta_2)^4) = 2/(\eta_1-\eta_2)^4$.

The factor of $1/3$ does \emph{not} come from the $\delta_{ij}$ contraction alone; it
comes from the angular average of $\hat{s}^i\hat{s}^j$ over all orientations:
%
\begin{equation}
  \langle\hat{s}^{i}\hat{s}^{j}\rangle = \frac{\delta^{ij}}{3} \,.
\end{equation}
%
This applies when the superposition orientation is not fixed but random, or when the
paper derives a result valid for a specific orientation and then states the formula for
the ``averaged'' case.  \textbf{FLAG FOR AUTHORS:} Clarify whether the $1/3$ is an
actual average over dipole orientations or whether it arises from the specific geometry
of the computation (e.g., taking only the component of the Wightman function relevant
to the physical decoherence).  The paper should state explicitly whether
$\Nnum$ is for a fixed orientation $\hat\mathbf{s}$ or an orientation-averaged result.

\section*{Summary of validity conditions}

\begin{center}
\begin{tabular}{lll}
\hline
Spacetime & Condition & Size of next correction \\
\hline
Power-law ($p>1$) & $d_{0}\ll u_{\rm on}/p$ & $(pd_0/u_{\rm on})^2$ \\
De Sitter & $d_{0}\ll 1/H$ & $(d_{0}H)^2$ \\
\hline
\end{tabular}
\end{center}

Both conditions are trivially satisfied for any laboratory-scale superposition in a
cosmological background ($d_0 H \lesssim 10^{-35}$ for $d_0\sim 1$~nm and $H\sim H_0$).
The dipole approximation is exact at the precision relevant for all physically realizable experiments.

\end{document}
